% ============================================================================
% RETRACTED -- DO NOT CITE, DO NOT SUBMIT.
%
% Every quantitative claim in this draft is an artifact. The fitness function
% could not see the property being measured: E_param = groups/P is exactly
% indifferent to placement, so a fully-occupied single-group genome satisfies
% the convolution predicate for free, and any operator that added units looked
% like a discovery. Controls: an operator copying NOTHING scored 26%, one that
% deliberately destroyed the target spacing scored 59%, the operator built to
% produce it scored 17%. A hardcoded, never-swept padd=0.006 was solely
% responsible for the 0/200 baseline.
%
% Kept only as a record of the failure. See FINDINGS.md and PROTOCOL.md.
% ============================================================================

\documentclass[11pt]{article}
\usepackage[a4paper,margin=1in]{geometry}
\usepackage[utf8]{inputenc}
\usepackage[T1]{fontenc}
\usepackage{amsmath,amssymb}
\usepackage{graphicx}
\usepackage{booktabs}
\usepackage[numbers]{natbib}
\usepackage{hyperref}

\title{\textbf{Copying Operators Make Repeated Structure Discoverable\\
in Evolutionary Architecture Search}}
\author{Vasili Gavrilov\thanks{Independent researcher. Reference implementation:
\url{https://github.com/Anode1/SMBPANN}.}}
\date{July 2026}

\begin{document}
\maketitle

\begin{center}\fbox{\parbox{0.9\textwidth}{\centering\textbf{RETRACTED.}
Every quantitative claim below is an artifact of an objective that could not see the
property being measured. Kept as a record of the failure only. See \texttt{FINDINGS.md}.}}\end{center}

\begin{abstract}
A companion study grew network topology from a dense seed under an evolutionary energy budget and
established four things: sparse task-relevant connectivity emerges on its own; weight sharing is adopted
when translation invariance rewards it; a compact filter emerges only once mutation acts on the whole
shared feature rather than on one edge; and compositional structure on the channel axis does not emerge
from a composite label alone. It left one piece imposed rather than emergent, because its genome was an
offset mask and so made the connectivity translation-invariant by construction: the \emph{tiling}, one
filter repeated across positions. This paper removes that imposition and supplies the missing operator.

We evolve topology from a disordered seed on a circular domain with nothing in the genome encoding
translation. Adding a single mutation that \emph{copies existing structure} raises the discovery rate of
a convolution from $0/200$ to $124/200$ ($p = 6.2\times10^{-16}$), with equivariance error falling from
$0.180$ to $0.046$ and accuracy rising toward the hand-built ceiling. The operative property is copying,
not scale: a local move that duplicates one unit at a time at its existing spacing also works
($34/200$, $p = 2.5\times10^{-11}$). The effect holds from $N=8$ to $24$, under two objective forms and
two pooling rules. We then show why such an operator is \emph{necessary} rather than merely convenient:
a planted convolution survives the per-site mutation that cannot build one ($122/150$ against
$0/200$, $p = 1.2\times10^{-15}$), and planting eleven of twelve units still leaves the last unreachable,
so the failure is neither instability nor insufficient search. A corollary for practice is that no
parameter-counting or connection-counting penalty can produce translational order, because neither has
order in its objective: the prior has to enter through the move set. We also contribute an
architecture-agnostic test for equivariance, and note that a topology search seeded from a
fully-occupied mask cannot support a claim that order emerged, because the seed supplies it.
\end{abstract}

\section{Introduction}

A companion study~\cite{paper1} used an evolutionary energy budget as a lens on convolutional structure,
growing topology from a dense seed and asking which pieces are imposed priors and which emerge. It
reported four results: (i)~sparse task-relevant connectivity emerges cleanly under a per-connection cost;
(ii)~weight sharing is adopted once translation invariance rewards it; (iii)~a compact aligned filter
emerges only when mutation acts on the whole shared feature rather than on a single edge, because sharing
decouples connection count from parameter count and a per-edge move therefore has no energy gradient;
and (iv)~compositional structure on the channel axis does not emerge from a composite label alone.

That study named its own limitation in its source: its genome was an offset mask, which ``makes the
connectivity translation-invariant by construction''. So the \emph{tiling}, one filter repeated across
positions, was imposed rather than emergent, and result (iii) concerned the filter alone. This paper
takes up that thread.

Decomposed, a convolution is three things: \textbf{locality}, a bounded receptive field; \textbf{sharing},
one filter used by many units, which is a statement about \emph{composition}; and \textbf{tiling}, those
units placed at regular intervals, which is a statement about \emph{order}. Locality is imposed here and
not studied (Section~\ref{sec:limits}). Composition, we find, needs no help at all: filters collapse onto
a single group under every cost term we tried, including one that charges nothing for sharing, driven by
data scarcity rather than by the budget. Order is the difficult one, and it is the subject of this paper.

\paragraph{The contribution.} A single additional mutation operator, one that \emph{copies existing
structure} rather than perturbing it at random, takes the discovery rate of a convolution from $0/200$ to
$124/200$. We show that the operative property is copying rather than scale, since a strictly local
version, duplicating one unit at a time at the spacing already present, also works. We then establish
that such an operator is necessary rather than convenient, by showing that a convolution is a stable
fixed point of the same dynamics that cannot construct it. We then draw the consequence for practice: no
parameter-counting or connection-counting penalty produces translational order, because neither measures
it, so the prior must enter through the move set. Two supporting contributions are an
architecture-agnostic test for equivariance, and a constraint on the seed: a search started from a
fully-occupied mask inherits the order it appears to discover.

This also revises how result (iii) of the companion study should be read. We had taken the per-offset
move to work because of its \emph{granularity}; on the present evidence it works because it copies one
decision across every instance of a shared feature, which is the same property identified here one level
up.

\section{Method}

\paragraph{Genome.} Over $P = N$ positions on a circle, \texttt{on[p]} says whether a hidden unit reads
the $K$-window starting at $p$, and \texttt{g[p]} names its weight group, with as many groups available
as positions. Within a group the filter, the bias and the readout weight are all shared, making a group
one channel; units in a group are max-pooled before the readout, which is what a convolutional network
does for detection. Nothing in the genome encodes translation.

\paragraph{The seed must be disordered.} A fully-occupied placement, of the kind a dense or
locally-connected seed provides, \emph{is} a perfect stride-1 tiling. Positional order is then present at
generation zero and the search need only merge filters, so such a seed measures whether order survives
rather than whether it forms. Under a fully-occupied seed and a low mutation rate the control reaches
$28\%$ convolutions with coverage and regularity both at $0.95$, which is the seed left undisturbed.
Every result below therefore starts from a \emph{disordered} seed: a random subset of positions with
random filters, so any order that appears has to be built. This constrains any topology search that
wants to claim order emerged.

\paragraph{Task.} A $K$-tap motif is planted at a uniformly random circular position, or is absent. A
unit sees the motif only if its window covers that position, so accuracy is bounded by coverage.
Training data is deliberately scarce ($24$ examples), following the companion study's finding that the
weight-sharing advantage grows as data thins.

\paragraph{Energy.} $E_{\text{param}} = \text{groups}/P$ pays per distinct weight, so sharing is dear and
placement free. $E_{\text{conn}} = \text{places}/P$ pays per connection, so placement is dear and sharing
free. Both are per-neuron fractions, so the budget scales with the network rather than penalising size.
Fitness is $(\text{acc} \geq \tau)\,?\,(2 - E) : \text{acc}$; Section~\ref{sec:robust} repeats everything
under a continuous $\text{acc} - \lambda E$.

\paragraph{Three mutation regimes.} \emph{Per-site} mutation flips occupancy at each position
independently and reassigns groups: the companion study's model, and our control. \emph{Segment repeat}
picks a period $T$ and repeats the genome's first $T$ slots across the range, a global rewrite, and
biologically the serial homology that makes a millipede. \emph{Tandem duplication} picks an occupied
position, measures the spacing to the nearest unit sharing its filter, and places one more unit of that
filter at that spacing: local, incremental, one unit at a time, with selection judging each addition.
The last is the crystal-growth move, and it is what decides whether scale or copying is the operative
property.

\paragraph{What counts as a convolution.} One shared filter, placements on a regular circular stride, and
no uncovered arc wider than the kernel. This admits a \emph{strided} convolution, which is still one.

\section{Equivariance, measured rather than inspected}\label{sec:equiv}

A placement diagram does not tell us whether what the search found ``is a convolution''. Theory does: a
linear map is translation-equivariant if and only if it is a convolution, stated for compact group
actions by Kondor and Trivedi~\cite{kondor}, with Cohen and Welling~\cite{cohenwelling} as predecessor.
We therefore test the property rather than inspect the wiring. A translation-equivariant detector fires
at the same rate wherever the motif sits, so we plant the motif at each position in turn and take the
spread of the detection rate; zero spread is equivariance. The measure assumes nothing about
architecture and transfers to any network.

The domain is circular because the theorem is stated for the cyclic group: on a finite strip a
convolution is genuinely not equivariant, since interior positions sit under more overlapping windows
than edge positions. Within a group the filter, the bias and the readout weight are all shared and the units are pooled, which
is what makes the composed map equivariant; sharing the filter alone leaves an unshared
locally-connected layer above it and the composition is not equivariant.

\section{Results}

\subsection{A copying operator makes the tiling discoverable}

From a disordered seed, per-site mutation collapses the filters onto a single group in nearly every run
but never assembles them into a regular placement: $0$ convolutions in $200$ runs. Adding one mutation
that repeats an existing segment of the genome raises this to $124$ of $200$
($p = 6.2\times10^{-16}$, Fisher exact, two-sided). Table~\ref{tab:main} gives the comparison and
Figure~\ref{fig:gallery} shows it run by run from shared starting points.

\begin{table}[t]
\centering\small
\begin{tabular}{@{}lccccc@{}}
\toprule
mutation regime & coverage & groups & regularity & test & \%conv \\
\midrule
\emph{convolution (ceiling)}   & $1.00$ & $1.0$ & $1.00$ & $0.918$ & n/a \\
\midrule
per-site only (control)        & $0.66$ & $1.1$ & $0.64$ & $0.843$ & $0\%$ \\
$+$ tandem duplication (local) & $0.86$ & $1.2$ & $0.85$ & $0.890$ & $17\%$ \\
$+$ segment repeat (global)    & $0.83$ & $1.0$ & $0.89$ & $0.885$ & $62\%$ \\
\midrule
random, matched budget         & $0.52$ & $3.2$ & $0.54$ & $0.700$ & $0\%$ \\
\bottomrule
\end{tabular}
\caption{$N=12$, $200$ seeds, disordered seed, max-pooling, $E_{\text{param}}$, per-site rate $0.002$.
Filters collapse onto one group in every regime, including the control. Only the copying regimes produce
regular placement. Fisher exact against the control: tandem duplication $p = 2.5\times10^{-11}$, segment
repeat $p = 6.2\times10^{-16}$.}
\label{tab:main}
\end{table}

The gain is not an artifact of how a convolution is scored. Equivariance error, which is continuous and
needs no threshold, falls from $0.180$ to $0.046$ over the same comparison, and test accuracy rises from
$0.843$ to $0.885$ against a hand-built ceiling of $0.918$. Scoring a convolution at relaxed regularity
thresholds of $0.90$ and $0.80$ rather than $0.999$ leaves the ordering unchanged.

\subsection{The operative property is copying, not scale}

The segment repeat rewrites the whole individual, so on its own it would suggest that a global move is
required. It is not. Tandem duplication takes one occupied position, measures the spacing to the nearest
unit sharing its filter, and places one further unit of that filter at that spacing. The move is local
and incremental, and selection judges each addition separately. It reaches $34$ of $200$
($p = 2.5\times10^{-11}$), against $0$ for per-site mutation at the same rate.

A move can therefore be as local as per-site mutation and still build order, provided it propagates what
is already present instead of randomising it. What separates the regimes is copying against perturbing,
not global against local. Figure~\ref{fig:structures} shows that the two operators also settle on
different lattices, a dense stride-1 tiling under tandem duplication and a strided one under the segment
repeat, both of which satisfy the definition. The operator's prior selects which convolution appears, not
only whether one does.

\subsection{Why the operator is necessary: order is stable but unreachable}\label{sec:reach}

An operator that helps might only be saving time. Running identical dynamics from two starting points
shows it is not. Plant a complete convolution in the seed and run per-site mutation alone: at
$\mu = 0.002$ it survives in $122$ of $150$ runs, with regularity $0.98$ and a mean ordered domain of
$11.7$ out of $12$. Start the same dynamics from disorder and it yields a convolution in $0$ of $200$
($p = 1.2\times10^{-15}$).

\begin{table}[t]
\centering\small
\begin{tabular}{@{}lccccc@{}}
\toprule
 & \multicolumn{3}{c}{planted, then per-site mutation} & \multicolumn{2}{c}{from disorder} \\
\cmidrule(lr){2-4}\cmidrule(lr){5-6}
$\mu$ & $12/12$ planted & $11/12$ planted & domain & \%conv & domain \\
\midrule
$0.030$ & $29\%$ & $2\%$ & $8.3$  & $1\%$ & $4.6$ \\
$0.008$ & $52\%$ & $3\%$ & $10.6$ & $0\%$ & $4.7$ \\
$0.002$ & $81\%$ & $6\%$ & $11.7$ & $0\%$ & $4.8$ \\
\bottomrule
\end{tabular}
\caption{A convolution is stable under the dynamics that cannot build it. ``Domain'' is the mean longest
run of consecutive units sharing one filter, out of $12$. Planted rows: $150$ seeds. Disordered rows:
$200$ seeds.}
\label{tab:stab}
\end{table}

Two alternative explanations are ruled out. The failure is not instability, since per-site mutation
preserves a convolution once it exists. And it is not a slow approach, since from disorder the ordered
domain grows to $4.8$ of $12$ and then stops: the same $4.8$ at $\mu = 0.002$ and at $\mu = 0.0002$,
a tenfold reduction in the disordering rate.

Planting $11$ of the $12$ units gives only $6\%$, against $81\%$ for a complete plant
($p = 6.8\times10^{-16}$). The search cannot add the twelfth unit. Closing the last gap requires one
specific position with one specific filter, and nothing in the objective points there, so the obstacle is
present at every step of assembly rather than only at the beginning. This is a reachability problem, and
placing order by construction is what the copying operators do about it.

\subsection{Dissolution sets whether order survives, not whether it can start}

Table~\ref{tab:ratio} varies the per-site rate, which is the rate at which existing structure is
destroyed, against fixed copying rates.

\begin{table}[t]
\centering\small
\begin{tabular}{@{}lccc@{}}
\toprule
per-site rate & control & $+$ tandem duplication & $+$ segment repeat \\
\midrule
$0.030$ & $1\%$ & $5\%$  & $21\%$ \\
$0.008$ & $0\%$ & $11\%$ & $40\%$ \\
$0.002$ & $0\%$ & $17\%$ & $62\%$ \\
\bottomrule
\end{tabular}
\caption{Order against the disordering rate, $200$ seeds each, disordered seed.}
\label{tab:ratio}
\end{table}

Both copying regimes improve as dissolution falls, which is the expected competition between adding and
removing units. The control stays at the floor at every rate. Reducing the disordering current therefore
governs whether order survives and does not affect whether it can start. Holding the copying rate fixed
and lowering dissolution reproduces the same effect, so the quantity that matters is the ratio rather
than either rate alone.

\subsection{The energy term selects composition, not order}

\begin{table}[t]
\centering\small
\begin{tabular}{@{}lccc@{}}
\toprule
mutation regime & $E_{\text{param}}$ & $E_{\text{conn}}$ & both \\
\midrule
per-site only (control)        & $0\%$  & $0\%$  & $1\%$ \\
$+$ tandem duplication (local) & $17\%$ & $0\%$  & $0\%$ \\
$+$ segment repeat (global)    & $62\%$ & $11\%$ & $10\%$ \\
\midrule
groups, all regimes            & $1.0$--$1.2$ & $1.1$--$1.3$ & $1.0$ \\
coverage, all regimes          & $0.66$--$0.86$ & $0.38$--$0.45$ & $0.38$--$0.44$ \\
\bottomrule
\end{tabular}
\caption{\%conv by cost term, $200$ seeds, disordered seed. Filters collapse onto one group under every
cost term, including $E_{\text{conn}}$, which charges nothing for a distinct weight.}
\label{tab:energy}
\end{table}

Composition needs no help from the budget. Filters collapse onto a single group under every cost term in
Table~\ref{tab:energy}, including $E_{\text{conn}}$, which offers no incentive to share at all. What
drives it there is data scarcity, since a shared filter learns from every position at once.

Order is selected by neither term. $E_{\text{param}}$ is indifferent to placement once filters are
shared, and $E_{\text{conn}}$ prevents order by thinning the network to roughly a third coverage, where
no regular stride survives. Tandem duplication is extinguished altogether under $E_{\text{conn}}$, which
removes each unit the operator adds.

The consequence for practice is that no parameter-counting or connection-counting penalty can produce
translational order, because neither has order in its objective. Order is not cheap and unreached under
these regularisers; it is not what they measure. This is why convolutional networks specify the tiling
rather than learn it, and it is not a limitation of compute. The prior enters through the architecture or
through the move set. Under search it does not disappear, it moves: in our case into the operators'
distributions over periods and lattice spacings, which were still chosen by hand.

\subsection{Robustness}\label{sec:robust}

Under a continuous objective $\text{acc} - \lambda E$, which has no threshold to go stale, the direction
of every result is unchanged at every $\lambda$ from $0.02$ to $1.60$, while the magnitude is smaller and
settles near $20\%$. The gap to the threshold form does not close as $\lambda$ grows. A threshold
objective is a constrained optimisation, in which energy is minimised with effectively unlimited weight
among individuals that clear the bar, whereas a linear penalty can never move fitness by more than
$\lambda$ against an accuracy range of roughly $0.4$. These are the constrained and penalty formulations
of the same problem and they are not interchangeable. A study using a threshold objective to elicit
structure relies on its lexicographic behaviour whether or not it says so.

Across problem size $N = 8$ to $24$ the control never exceeds $4\%$ while the copying regimes always
exceed it significantly. Under sum-pooling rather than max-pooling the magnitudes are larger, and the reference is no longer a
ceiling because sum-pooling accumulates noise from every position while the motif occupies one; we report
max-pooling as primary for that reason.

\begin{figure}[t]
\centering
\includegraphics[width=\textwidth]{fig_structures.pdf}
\caption{The structures as weight matrices: rows are hidden units, columns are input positions, a cell is
filled when that unit reads that input, and colour is the weight group, so the same colour is the same
weight. Sharing is colour uniformity and order is a regular band. A convolution is a circulant matrix,
one colour and constant along each diagonal, wrapping at the corner. Per-site mutation collapses the
colour to one but does not close the band. The two copying operators close it on different lattices,
dense in one case and strided in the other, and both satisfy the definition.}
\label{fig:structures}
\end{figure}

\begin{figure}[t]
\centering
\includegraphics[width=\textwidth]{fig_gallery.pdf}
\caption{Eight disordered starting points put through each move set, sampled at a regular stride across
the $200$ runs rather than selected. Columns share a seed, so differences down a column are due to the
move set alone. Colour collapses to one filter in nearly every panel of every row, including the control,
so composition is not the difficult part. What separates the rows is whether the band closes.}
\label{fig:gallery}
\end{figure}

\section{Limitations}\label{sec:limits}

Locality is imposed rather than searched: the genome offers only $K$-windows, so unlike the companion
study this work cannot ask whether a bounded receptive field arises on its own, and addresses only what
happens to composition and order once locality is granted. The tasks are small synthetic constructions
built so that the intended structure is optimal. Subsampling, the second half of the Neocognitron's
architecture, is absent entirely. The copying operators carry their own priors, in the distributions over
periods and lattice spacings, so they relocate the commitment rather than removing it. We give no quantitative account of the domain size at which
per-site dynamics stops, only the measurement that it stops at roughly $4.8$ of $12$ and does not move
when the mutation rate is cut tenfold.

\section{Related work}

The prior this paper decomposes was not itself found by search. Hubel and Wiesel~\cite{hubelwiesel} read
local receptive fields, simple and complex cells, and orientation columns off the visual cortex; the
Neocognitron~\cite{fukushima} transplanted that architecture into an artificial network as its S-cell and
C-cell layers; explicit weight sharing and backpropagation training followed~\cite{lecun89}. Convolution
entered machine learning as a borrowed biological prior rather than an emergent one.

Convolutional structure has been shown to emerge in fully-connected networks trained by gradient
descent~\cite{ingrosso} and learned from scratch by regularisation~\cite{neyshabur}. Those results obtain
locality and translation structure without imposing them, by a gradient route over weights rather than a
search over topology, and the present negative does not contradict them: a gradient on a continuous
weight space is not a per-site move on a discrete placement. Using an evolutionary connection cost as a
lens on emergent structure has precedent, with a connection cost driving modularity~\cite{clune} and
hierarchy~\cite{mengistu}, and Weight Agnostic Networks~\cite{wann} isolating what topology alone
contributes. What is new here is the separation of composition from order, the finding that only the
first is reachable by undirected search, and the identification of copying as the operative property of
the moves that reach the second.

\begin{thebibliography}{9}\setlength{\itemsep}{0pt}
\bibitem{paper1} V.~Gavrilov. The Imposed and Emergent Pieces of Convolution Under an Energy Budget.
 2026. \url{https://doi.org/10.5281/zenodo.21423177}.
\bibitem{hubelwiesel} D.~H. Hubel, T.~N. Wiesel. Receptive fields, binocular interaction and functional
 architecture in the cat's visual cortex. \emph{Journal of Physiology} 160(1), 1962.
\bibitem{fukushima} K.~Fukushima. Neocognitron: A self-organizing neural network model for a mechanism
 of pattern recognition unaffected by shift in position. \emph{Biological Cybernetics} 36(4), 1980.
\bibitem{lecun89} Y.~LeCun et al. Backpropagation Applied to Handwritten Zip Code Recognition.
 \emph{Neural Computation} 1(4), 1989.
\bibitem{kondor} R.~Kondor, S.~Trivedi. On the Generalization of Equivariance and Convolution in Neural
 Networks to the Action of Compact Groups. \emph{ICML} 2018.
\bibitem{cohenwelling} T.~S. Cohen, M.~Welling. Group Equivariant Convolutional Networks.
 \emph{ICML} 2016.
\bibitem{ingrosso} A.~Ingrosso, S.~Goldt. Data-driven emergence of convolutional structure in neural
 networks. \emph{PNAS} 119(40), 2022.
\bibitem{neyshabur} B.~Neyshabur. Towards Learning Convolutions from Scratch. \emph{NeurIPS} 2020.
\bibitem{clune} J.~Clune, J.-B. Mouret, H.~Lipson. The evolutionary origins of modularity.
 \emph{Proc. R. Soc. B} 280(1755), 2013.
\bibitem{mengistu} H.~Mengistu, J.~Huizinga, J.-B. Mouret, J.~Clune. The Evolutionary Origins of
 Hierarchy. \emph{PLOS Comput. Biol.} 12(6), 2016.
\bibitem{wann} A.~Gaier, D.~Ha. Weight Agnostic Neural Networks. \emph{NeurIPS} 2019.
\end{thebibliography}

\end{document}
